\documentclass{amsart}
\usepackage{amssymb,mathtools,array,longtable,seqsplit}
\usepackage[hidelinks]{hyperref}
\usepackage{url}
\hypersetup{pdftitle={Irrationality of the dilogarithm at 1/2},pdfauthor={Ruiyi Zhang}}
\setlength{\emergencystretch}{3em}
\newtheorem{theorem}{Theorem}
\newtheorem{lemma}{Lemma}[section]
\newtheorem{proposition}[lemma]{Proposition}
\numberwithin{equation}{section}
\numberwithin{table}{section}
\newcommand{\Li}{\operatorname{Li}}
\newcommand{\Res}{\operatorname*{Res}}
\newcommand{\Z}{\mathbb Z}
\newcommand{\fl}[1]{\lfloor #1\rfloor}
\newcommand{\br}[1]{[\,#1\,]}

\begin{document}
\title{Irrationality of the dilogarithm at 1/2}
\author{Ruiyi Zhang}
\address{Independent Researcher, Columbia, Maryland, USA}
\email{ryzfsu@gmail.com}
\subjclass[2020]{Primary 11J72; Secondary 33B30}
\keywords{dilogarithm, irrationality, determinant polynomials, logarithmic energy, interval certificates, prime-power congruences}
\begin{abstract}
We prove that $\operatorname{Li}_2(1/2)$ is irrational. Irrationality of
$\operatorname{Li}_2(1/s)$ was known for integers $s\le-5$ and $s\ge6$ by
work of Hata and of Rhin and Viola, and Prellberg recently proved it for
$s=-4,5,-3$. The approximations come from determinants of moments of
$1/(xy-2)$, differentiated $m/2$ times in the parameter. We bound their
values by a logarithmic energy on an ellipse, certified by directed interval
arithmetic. Valuations of binomial coefficients give a common denominator of
the coefficients. A Lucas congruence modulo $83$ gives nonvanishing along
$m=32\cdot83^k$. Three finite computations are part of the proof.
\end{abstract}
\maketitle

\section{Introduction}\label{sec:introduction}

For $|w|\le1$ let $\operatorname{Li}_2(w)=\sum_{n\ge1}w^n/n^2$. Hata
\cite[Theorem~1.1]{Ha} proved that $\operatorname{Li}_2(1/s)$ is irrational
for every integer $s\le-5$ and every $s\ge7$. Rhin and Viola adapted their
permutation group method \cite{RV96} to the dilogarithm and settled $s=6$
\cite{RV05}. They later treated negative arguments \cite{RV19}. Viola and
Zudilin \cite{VZ} introduced a six-parameter family of double integrals. In
2024 Calegari, Dimitrov and Tang \cite[Remark~2.8.2 and Section~15.2]{CDT}
listed the integers still open: $s\in\{-4,-3,-2,2,3,4,5\}$. Prellberg's 2026
preprint \cite{Pr} treats $s=-4,5,-3$. We prove the case $s=2$.

Put
\[
 \beta=\operatorname{Li}_2(1/2)=\sum_{r=1}^{\infty}\frac{2^{-r}}{r^2},
 \qquad \ell=\log 2.
\]
For each positive integer $m$ divisible by $4$, let
\[
 g=\frac m2,\qquad
 I_m=\left[\frac m4,\frac{5m}4\right)\cap\mathbb Z,
 \qquad J_m=[0,2m)\cap\mathbb Z.
\]
The rational matrices $A^{(g)},B^{(g)},C^{(g)}$ are defined in
Section~\ref{sec:source}. Rows and columns are in increasing order. Set
\begin{equation}\label{eq:main-polynomial}
 G_m(T)=\det\begin{bmatrix}
 C^{(g)}_{I_m,J_m}\\
 B^{(g)}_{I_m,J_m}-T A^{(g)}_{I_m,J_m}
 \end{bmatrix}\in\mathbb Q[T].
\end{equation}
Only the last $m$ rows depend on $T$, so $\deg G_m\le m$.

\begin{theorem}\label{thm:main}
The number $\operatorname{Li}_2(1/2)$ is irrational.
\end{theorem}

By the reflection formula \cite[\S25.12]{DLMF},
$\operatorname{Li}_2(1/2)=\pi^2/12-\frac12\log^22$, so Theorem~\ref{thm:main} is
equivalent to the irrationality of $\pi^2-6\log^22$.

Theorem~\ref{thm:main} follows from three estimates for the polynomials
\eqref{eq:main-polynomial}.

\begin{proposition}\label{prop:main}
For every $m\ge8$ with $4\mid m$,
\begin{equation}\label{eq:analytic}
 |G_m(\beta)|\le
 \exp\left(-\alpha m^2+10m\log m+56m+301\right),
 \qquad \alpha=\frac{1992909}{10^6}.
\end{equation}
For every positive $m$ divisible by $4$ there is a positive rational number
 $D_m$ such that
\begin{equation}\label{eq:integerizer}
 D_mG_m(T)\in\mathbb Z[T],\qquad
 \limsup_{\substack{m\to\infty\\4\mid m}}\frac{\log D_m}{m^2}
 \le \tau<1.979.
\end{equation}
For every fixed $t\in\mathbb Q$,
\begin{equation}\label{eq:nonvanishing}
 G_{32\cdot83^k}(t)\ne0
 \quad\text{if }k\ge0\text{ and }2k+v_{83}(t)>0,
\end{equation}
where $v_{83}(0)=+\infty$.
\end{proposition}

Section~\ref{sec:source} writes $G_m(\beta)$ as a determinant of rational
series. Sections~\ref{sec:contour}--\ref{sec:frequency} prove
\eqref{eq:analytic}. Section~\ref{sec:integerization} proves
\eqref{eq:integerizer}, and Section~\ref{sec:nonvanishing} proves
\eqref{eq:nonvanishing}. Theorem~\ref{thm:main} follows in
Section~\ref{sec:proof}.

Since $1.979-1.992909<0$, Theorem~\ref{thm:main} follows.

Three finite computations are part of the proof (Appendix~\ref{app:ver}).
Directed interval inequalities give the constant $\alpha$. Exact integer
arithmetic gives the coefficient bounds. Modular arithmetic gives the seed
congruence. The use of AI tools is described at the end of the paper.

\section{Rational matrices and parameter derivatives}\label{sec:source}

Fix an integer $g\ge0$. All indices $i,j$ are nonnegative integers. Empty
sums are zero. Put
\[
U_i=2^i\binom{i+g}{g},\qquad
t_i=\sum_{r=1}^{i}\frac{U_{i-r}}r
 +2^i\sum_{s=1}^{g}\frac{(-1)^{s-1}}s
       \binom{i+g}{g-s}(2^s-1).
\]
For $i\ne j$ define
\[
C^{(g)}_{ij}=\frac{U_i-U_j}{j-i},\qquad
B^{(g)}_{ij}=\frac{t_i-t_j}{j-i},\qquad
A^{(g)}_{ij}=0.
\]
The diagonal entries are
\begin{equation}\label{eq:source-diag-AC}
A^{(g)}_{ii}=U_i,\qquad
C^{(g)}_{ii}=-U_i\sum_{s=1}^{g}\frac1{i+s},
\end{equation}
and
\begin{equation}\label{eq:source-diag-B}
B^{(g)}_{ii}
 =\sum_{r=1}^{i}\frac{U_{i-r}}{r^2}
 -2^i\sum_{s=2}^{g}\frac{(-1)^s}{s}\binom{i+g}{g-s}
       \sum_{r=1}^{s-1}\frac{2^r-1}{r}.
\end{equation}
These are rational matrices. Their diagonals are defined separately from
the off-diagonal quotients.

For real $a>1$, let
\[
\begin{aligned}
t_i^0(a)&=\sum_{r=1}^{i}\frac{a^{i-r}}r,&
b_i^0(a)&=\sum_{r=1}^{i}\frac{a^{i-r}}{r^2},\\
L(a)&=\log(1-1/a)
     =-\sum_{r=1}^{\infty}\frac{a^{-r}}r,&
\beta(a)&=\operatorname{Li}_2(1/a)
     =\sum_{r=1}^{\infty}\frac{a^{-r}}{r^2}.
\end{aligned}
\]
The logarithm is real. Define
\[
A^0_{ij}(a)=a^i\delta_{ij}.
\]
\[
C^0_{ij}(a)=
\begin{cases}
\dfrac{a^i-a^j}{j-i},&i\ne j,\\[4pt]
0,&i=j,
\end{cases}
\]
\[
B^0_{ij}(a)=
\begin{cases}
\dfrac{t_i^0(a)-t_j^0(a)}{j-i},&i\ne j,\\[4pt]
b_i^0(a),&i=j,
\end{cases}
\]
where $\delta_{ij}$ is the Kronecker delta.

\begin{lemma}\label{lem:source-mellin}
For $a>1$ and $i,j\ge0$,
\begin{equation}\label{eq:source-mellin}
\begin{aligned}
M_{ij}(a)
&:=B^0_{ij}(a)+L(a)C^0_{ij}(a)-\beta(a)A^0_{ij}(a)\\
&=-\sum_{r=1}^{\infty}\frac{a^{-r}}{(i+r)(j+r)}.
\end{aligned}
\end{equation}
Equivalently,
\[
M_{ij}(a)=-\int_0^1\!\!\int_0^1
                  \frac{x^i y^j}{a-xy}\,dy\,dx.
\]
\end{lemma}

\begin{proof}
Split off the first $i$ terms of the defining series. Then
\[
t_i^0(a)+a^iL(a)
 =-\sum_{r=1}^{\infty}\frac{a^{-r}}{i+r},\qquad
b_i^0(a)-a^i\beta(a)
 =-\sum_{r=1}^{\infty}\frac{a^{-r}}{(i+r)^2}.
\]
For $i\ne j$, take the difference of the first identity at $i$ and at $j$,
and divide by $j-i$. By
\[
\frac1{j-i}\left(\frac1{i+r}-\frac1{j+r}\right)
 =\frac1{(i+r)(j+r)},
\]
this proves \eqref{eq:source-mellin} off the diagonal. The second identity
proves it on the diagonal, where $C^0_{ii}=0$.

Finally,
\[
\frac1{a-xy}=\sum_{r=1}^{\infty}a^{-r}(xy)^{r-1}.
\]
For fixed $a>1$ this series converges uniformly and absolutely on
$[0,1]^2$. Termwise integration gives the integral formula.
\end{proof}

Put
\[
D_g f(a)=\frac1{g!}\frac{\partial^g}{\partial a^g}
                         \bigl(a^g f(a)\bigr).
\]
For every integer $r\ge1$,
\[
D_g a^{-r}
 =\frac1{g!}\prod_{h=0}^{g-1}(g-r-h)\,a^{-r}
 =
\begin{cases}
0,&1\le r\le g,\\[2pt]
(-1)^g\binom{r-1}{g}a^{-r},&r>g.
\end{cases}
\]
This is the $g$th derivative of $a^{g-r}$ divided by $g!$. For
$1\le r\le g$ the polynomial $a^{g-r}$ has degree less than $g$. For $r>g$
the $g$ negative factors give the binomial coefficient. For $g=0$ the
product is empty and equals one.

Fix a compact subinterval of $(1,\infty)$. On it, the derivatives through
order $g$ of the summands in $a^gM_{ij}(a)$ are bounded by $C(1+r)^gq^r$.
Here $C>0$ and $0<q<1$ do not depend on $r$. Hence
\eqref{eq:source-mellin} may be differentiated termwise, and
\[
(D_gM_{ij})(a)
 =(-1)^{g+1}\sum_{r>g}
       \frac{\binom{r-1}{g}a^{-r}}{(i+r)(j+r)}.
\]
Put
\begin{equation}\label{eq:source-N}
N_{ij}:=(D_gM_{ij})(2)
 =(-1)^{g+1}\sum_{r>g}\frac{w_r}{(i+r)(j+r)},
\qquad
w_r=\binom{r-1}{g}2^{-r}.
\end{equation}
The weights are positive for $r>g$, and the series converges absolutely.

Next we find the rational coefficients of $D_gM$. For $a>1$ put
\[
U_i(a)=a^i\binom{i+g}{g},\qquad
Q_s(a)=\left(\frac{a}{a-1}\right)^s-1\quad(s\ge1).
\]
Define $t_i^{(g)}(a)$, $A^{(g)}_{ij}(a)$, $B^{(g)}_{ij}(a)$ and
$C^{(g)}_{ij}(a)$ by the same finite formulas. In them replace $U_n$ by
$U_n(a)$, $2^n$ by $a^n$, and each factor $2^s-1$ by $Q_s(a)$. Thus
\[
t_i^{(g)}(a)
 =\sum_{r=1}^{i}\frac{U_{i-r}(a)}r
  +a^i\sum_{s=1}^{g}\frac{(-1)^{s-1}}s
          \binom{i+g}{g-s}Q_s(a).
\]
These are rational functions of $a$, regular for $a>1$. At $a=2$ they equal
the rational quantities defined above.

\begin{lemma}\label{lem:source-bank}
For every $g\ge0$, $i,j\ge0$ and $a>1$,
\[
(D_gM_{ij})(a)
 =B^{(g)}_{ij}(a)+L(a)C^{(g)}_{ij}(a)
                   -\beta(a)A^{(g)}_{ij}(a).
\]
\end{lemma}

\begin{proof}
For every $C^g$ function $f$, Leibniz's rule gives
\[
D_g(a^if(a))
 =\sum_{s=0}^{g}\binom{i+g}{g-s}
                  a^{i+s}\frac{f^{(s)}(a)}{s!}.
\]
In particular $D_ga^i=U_i(a)$. As $L(a)=\log(a-1)-\log a$, for $s\ge1$
\[
\frac{a^sL^{(s)}(a)}{s!}
 =\frac{(-1)^{s-1}}sQ_s(a).
\]
Differentiating the absolutely convergent series for $\beta$ gives
\[
\beta'(a)
 =-\frac1a\sum_{r=1}^{\infty}\frac{a^{-r}}r
 =\frac{L(a)}a.
\]
Leibniz's rule for the $(s-1)$st derivative of $L(a)/a$ gives
\[
\begin{aligned}
\frac{a^s\beta^{(s)}(a)}{s!}
 &=\frac1s\sum_{r=0}^{s-1}
       (-1)^{s-1-r}\frac{a^rL^{(r)}(a)}{r!}\\
 &=\frac{(-1)^{s-1}}s
       \left(L(a)-\sum_{r=1}^{s-1}\frac{Q_r(a)}r\right).
\end{aligned}
\]

Hence
\[
D_g\bigl(t_i^0(a)+a^iL(a)\bigr)
 =t_i^{(g)}(a)+L(a)U_i(a).
\]
Divided differences in $i,j$ prove the lemma for $i\ne j$, including
$A^{(g)}_{ij}(a)=0$.

On the diagonal $M_{ii}=b_i^0-a^i\beta$, and the same calculation gives
\[
\begin{aligned}
(D_gM_{ii})(a)
={}&\sum_{r=1}^{i}\frac{U_{i-r}(a)}{r^2}
       -U_i(a)\beta(a)\\
 &-a^iL(a)\sum_{s=1}^{g}
       \frac{(-1)^{s-1}}s\binom{i+g}{g-s}\\
 &-a^i\sum_{s=2}^{g}\frac{(-1)^s}{s}\binom{i+g}{g-s}
            \sum_{r=1}^{s-1}\frac{Q_r(a)}r.
\end{aligned}
\]
The first and last sums give $B^{(g)}_{ii}(a)$, with the sign in
\eqref{eq:source-diag-B}. For the logarithmic coefficient we use
\[
\sum_{s=1}^{g}\frac{(-1)^{s-1}}s\binom{i+g}{g-s}
 =\binom{i+g}{g}\sum_{s=1}^{g}\frac1{i+s}.
\]
For its proof, differentiate at $x=0$ the polynomial identity
\[
\binom{i+g+x}{g}
 =\sum_{s=0}^{g}\binom{i+g}{g-s}\binom{x}{s}.
\]
It holds for nonnegative integers $x$ by the binomial theorem, hence as a
polynomial identity. By the product formula, the derivative of its left
side at zero is $\binom{i+g}{g}\sum_{s=1}^{g}(i+s)^{-1}$. For $s\ge1$, the
derivative of $\binom{x}{s}$ at zero is $(-1)^{s-1}/s$. This proves the
equality, also when $g=0$.

So the logarithmic coefficient is
$-U_i(a)\sum_{s=1}^{g}(i+s)^{-1}=C^{(g)}_{ii}(a)$. The coefficient of
$-\beta(a)$ is $U_i(a)=A^{(g)}_{ii}(a)$. This proves the diagonal case.
\end{proof}

Now let $m$ be a positive integer divisible by four. Set
\[
g=m/2,\qquad
I_m=\{m/4,\ldots,5m/4-1\},\qquad
J_m=\{0,\ldots,2m-1\}.
\]
Rows and columns are in increasing index order. The polynomial
\eqref{eq:main-polynomial} is
\[
G_m(T)=
\det\begin{bmatrix}
C^{(g)}_{I_m,J_m}\\
B^{(g)}_{I_m,J_m}-T A^{(g)}_{I_m,J_m}
\end{bmatrix}.
\]
It lies in $\mathbb Q[T]$. Only the last $m$ rows depend on $T$, each
affinely, so its degree is at most $m$.

Put $\beta=\beta(2)=\operatorname{Li}_2(1/2)$ and $\ell=\log2$.
As $L(2)=-\ell$, Lemma~\ref{lem:source-bank} gives
\[
N=B^{(g)}-\ell C^{(g)}-\beta A^{(g)}.
\]
In the determinant for $G_m(\beta)$, add $-\ell$ times each top row to the
corresponding bottom row. This is left multiplication by
\[
\begin{bmatrix}
\mathrm{Id}_m&0\\
-\ell\,\mathrm{Id}_m&\mathrm{Id}_m
\end{bmatrix},
\]
which has determinant one. Hence
\begin{equation}\label{eq:source-determinant}
G_m(\beta)=
\det\begin{bmatrix}
C^{(g)}_{I_m,J_m}\\
N_{I_m,J_m}
\end{bmatrix}.
\end{equation}
This row operation needs no rank or nonvanishing assumption.

\section{Exact contour expansion and constrained energy}\label{sec:contour}

Set
\[
 \Gamma=\{\xi+A\cos u+iB\sin u:0\le u<2\pi\},\quad
 (\xi,A,B)=\left(\frac{11281}{10000},
                \frac{12783}{10000},\frac{229}{200}\right).
\]
It encloses \([0,2]\). Before scaling, the contour in this section is
\(m\Gamma\). The entire function \(F_g(z)=2^z\binom{z+g}{g}\) satisfies
\begin{equation}\label{eq:contour-C}
 C^{(g)}_{ij}=-\oint_{m\Gamma}
       \frac{F_g(z)}{(z-i)(z-j)}\frac{dz}{2\pi i}
        +\ell U_i[i=j].
\end{equation}
For distinct poles this is the residue formula for divided differences. At
the double pole \(F_g'(i)=U_i(\ell+\sum_{s=1}^g1/(i+s))\), which gives
\eqref{eq:contour-C} on the diagonal too. No sine factor occurs.

Write \(K(x,y)=1/(x+y)\) and \(\Delta(X)=\prod_{a<b}(x_b-x_a)\) for
increasing real lists. The same product defines \(\Delta\) on complex lists.
Cauchy--Binet applied to \eqref{eq:source-N} gives the absolutely
convergent expansion
\begin{equation}\label{eq:contour-CB}
 G_m(\beta)=(-1)^{m(g+1)}
 \sum_{\substack{R\subset\{g+1,g+2,\ldots\}\\|R|=m}}
  \Bigl(\prod_{r\in R}w_r\Bigr)
 \det K_{I_m,R}\det\begin{bmatrix}C^{(g)}\\K_{R,J_m}\end{bmatrix}.
\end{equation}
The sum runs over subsets, so no \(m!\) occurs. For fixed \(m\) the weights
\(w_r\) decay exponentially and both minors are bounded. Finite truncation
and passage to the limit justify the expansion.

Select the diagonal term of \eqref{eq:contour-C} in the rows \(S\subset
I_m\). Put \(k=|S|\), \(d=m-k\), and let \(Z=(z_1,\ldots,z_d)\) lie on
\(m\Gamma\). The selected columns are exactly \(S\). Andr\'eief's
identity for the remaining rows gives the absolute integrand
\begin{equation}\label{eq:contour-Q}
 \begin{aligned}
 Q_m(S,R,Z)={}&\Bigl(\prod_{i\in S}\ell U_i\Bigr)
  \Bigl(\prod_{r\in R}w_r\Bigr)\Bigl(\prod_{z\in Z}|F_g(z)|\Bigr)\\
 &\times\left|\det K_{I_m\setminus S,-Z}\,\det K_{I_m,R}
               \,\det K_{(-Z,R),J_m\setminus S}\right|.
 \end{aligned}
\end{equation}
The measure is \(\prod|dz|/(2\pi)\) and the symmetrization factor is
\(1/d!\). By the triangle inequality \(|G_m(\beta)|\le E_m(\infty)\), where
\begin{equation}\label{eq:contour-E}
 E_m(H)=\sum_{S\subset I_m}\frac1{d!}
   \int_{(m\Gamma)^d}
    \sum_{\substack{|R|=m\\g<r\le H}}Q_m(S,R,Z)
                         \prod_{z\in Z}\frac{|dz|}{2\pi}.
\end{equation}

The signed product formula fixes every quadratic-scale normalization. Put
\(n=m/4\), \(Y=S\uplus Z\),
\[
 D_m(a)=(a-n)!(n+m-1-a)!a!(2m-1-a)!,
\]
\[
 P_m(z)=\prod_{a\in I_m\uplus J_m}(z-a),\qquad
 Q_m^*(r)=\prod_{a\in I_m\uplus J_m}(r+a).
\]
Then
\begin{equation}\label{eq:contour-signed}
\begin{split}
 G_m(\beta)={}&\Delta(I_m)^2\Delta(J_m)
 \sum_{S\subset I_m}\frac{(-1)^{k(n+1)}\ell^k}{d!}
 \sum_R\prod_{a\in S}\frac{U_a}{D_m(a)}\prod_{r\in R}w_r\\
 &\quad\times\oint_{(m\Gamma)^d}
 \frac{\Delta(Y)^2\Delta(R)^2\prod_{y\in Y,r\in R}(y+r)}
      {\prod_{z\in Z}P_m(z)\prod_{r\in R}Q_m^*(r)}
                  \prod_{z\in Z}F_g(z)\frac{dz}{2\pi i}.
\end{split}
\end{equation}
In this formula the three Cauchy determinants give the squared
Vandermondes and the single frequency cross product. The conversion
\[
 \Delta(I_m\setminus S)\Delta(J_m\setminus S)
 =\frac{\Delta(I_m)\Delta(J_m)\Delta(S)^2}
              {\prod_{a\in S}D_m(a)}
\]
restores the removed fixed factors. Each pinned index occurs twice in a
contour denominator; restoring both gives \((z-a)^2\). Restoring its one
removed frequency factor gives \((r+a)\). Since \(m\) is even, the Laplace
sign, the minus signs from \eqref{eq:contour-C} and the Cauchy signs reduce
to \((-1)^{k(n+1)}\). The upper estimate does not use the signs.

Scale all particle positions by \(m\). Write \(E=[1/4,5/4]\), and define
\[
 \begin{aligned}
 A_0(z)&=\int_E\log|z-x|\,dx+\int_0^2\log|z-x|\,dx,\\
 B_0(t)&=\int_E\log(t+x)\,dx+\int_0^2\log(t+x)\,dx.
 \end{aligned}
\]
\[
 \phi(z)=\Re(z)\ell+\int_0^{1/2}\log|z+s|\,ds
                         +\tfrac12-\tfrac12\log(1/2),
\]
\begin{equation}\label{eq:contour-psi}
 \psi(t)=t\log t-(t-1/2)\log(t-1/2)
                           -\tfrac12\log(1/2)-t\ell.
\end{equation}
The integral defines \(\phi\) also across the negative real axis, with
\(0\log 0=0\). So \(\phi\) has no branch ambiguity where \(\Gamma\) crosses
the interior of the zero segment \([-1/2,0]\).

Let \(K_0(z,w)=-\log|z-w|\), and write
\[
 I(\sigma,\tau)=\iint K_0(z,w)\,d\sigma(z)\,d\tau(w),
 \qquad I(\sigma)=I(\sigma,\sigma).
\]
For a measure \(\mu\) on the positive real line, let
\(\mu^-=(-\mathrm{id})_\#\mu\) be its reflection onto the negative real line.
The admissible measures are those with
\begin{equation}\label{eq:contour-admissible}
 \begin{aligned}
 0&\le\kappa\le1_E\,dx,\quad \nu\ge0,\quad
   \operatorname{supp}\nu\subset\Gamma,\quad |\kappa|+|\nu|=1,\\
 0&\le\mu\le1_{[1/2,\infty)}\,dt,\quad |\mu|=1.
 \end{aligned}
\end{equation}
The contour density has no cap, and \(|\kappa|\) is not prescribed.
With \(\eta=\kappa+\nu\), the energy is
\begin{equation}\label{eq:contour-energy}
 \begin{aligned}
 \mathcal E(\eta,\mu)={}&C_0+\int(A_0-\phi)\,d\eta
            +\int(B_0-\psi)\,d\mu\\
            &+I(\eta)+I(\mu)+I(\eta,\mu^-),\\
 C_0={}&\tfrac92-2\log2.
 \end{aligned}
\end{equation}
The pinned--contour interaction has coefficient two. Each
companion--frequency interaction has coefficient one. These coefficients
come from \eqref{eq:contour-signed}. For the fixed Vandermonde,
\[
 \log(\Delta(I_m)^2\Delta(J_m))
 =3m^2\log m+(2\log2-9/2)m^2+O(m\log m).
\]
Here \(\int_0^b\!\int_0^b\log|x-y|\,dx\,dy=b^2\log b-3b^2/2\). Twice the
half-energy for length one plus the half-energy for length two gives the
constant. Every \(m^2\log m\) term in \eqref{eq:contour-signed} cancels, for
every pinned mass. The remaining homogeneous factor is \(m^{-2m-2d}\)
before the contour differentials and \(m^{-2m-d}\) with them.

To define \eqref{eq:contour-energy} without \(\infty-\infty\), put
\(\ell_1(t)=\log(1+t)\). Write its frequency-dependent part as
\begin{equation}\label{eq:contour-regularized}
 \int(B_0-\psi-3\ell_1)\,d\mu
 +\iint\log\frac{(1+t)(1+s)}{|t-s|}\,d\mu(t)d\mu(s)
 +\iint\log\frac{1+t}{|z+t|}\,d\eta(z)d\mu(t).
\end{equation}
The self kernel is nonnegative. The mixed kernel is bounded on the compact
companion support times the half-line. The first field is
\(t\log2-\frac12\log t+O(1)\). Hence an infinite first moment gives energy
\(+\infty\). The same definition works on a sufficiently small tube about
\(\Gamma\). Define
\begin{equation}\label{eq:contour-alpha}
 \alpha_\Gamma=\inf\{\mathcal E(\kappa+\nu,\mu):
   \eqref{eq:contour-admissible}\}.
\end{equation}
This is the variational problem on the full half-line. An admissible trial
measure gives only an upper bound for this infimum. The lower bound comes
from the dual certificate in Section~\ref{sec:dual}.

\section{Circle regularization and transfer from the tube}\label{sec:tube}

Let \(T_\delta=\{z:\operatorname{dist}(z,\Gamma)\le\delta\}\). In
\eqref{eq:contour-admissible} replace only the support condition on \(\nu\)
by \(\operatorname{supp}\nu\subset T_\delta\). For \(0<\delta\le1/64\) we
have the exact estimate
\begin{equation}\label{eq:tube-transfer}
 \inf_{T_\delta}\mathcal E\ge\alpha_\Gamma-300\delta.
\end{equation}

The proof covers the unrestricted contour self-energy. Let \(D\) be the
interior of the ellipse. For a point \(z\) inside \(D\), let \(\omega_z\) be
its interior harmonic measure on \(\Gamma\). Outside \(D\), use exterior
harmonic measure with total mass one. At a boundary point use
\(\delta_z\). Set
\[
 c=(A+B)/2,\quad d=(A-B)/2,\quad
 F(w)=\xi+cw+d/w\quad(|w|\ge1).
\]
The factorization \(F(w_1)-F(w_2)=(w_1-w_2)(c-d/(w_1w_2))\) shows that
\(F\) is one-to-one outside the unit disk. Let \(W\) be its inverse. Set
\(e(z)=0\) inside the ellipse and \(e(z)=\log|W(z)|\) outside. The
Green-function identity gives
\begin{equation}\label{eq:tube-green}
 U^{\omega_z}(x)\le K_0(z,x)+e(z),\qquad
 I(\omega_z,\omega_w)\le K_0(z,w)+e(z)+e(w).
\end{equation}
For two interior points the first deficit is the interior Green function.
For points on opposite sides it is zero. For two exterior points it is the
nonnegative Green function of the exterior of the disk,
\[
 \log\left|\frac{W(z)\overline{W(x)}-1}{W(z)-W(x)}\right|.
\]
Its sign follows from
\(|a\overline b-1|^2-|a-b|^2=(|a|^2-1)(|b|^2-1)\ge0\). At boundary points
the identity in \eqref{eq:tube-green} is an equality. Integrating it
against \(\omega_w\) proves the two-sweep inequality. These identities also
follow by Dirichlet uniqueness after subtracting the Green pole, including
the growth term at infinity. Integrable logarithms and limiting arguments
cover boundary points and arbitrary finite-energy measures.

If \(z\in T_\delta\) is exterior, the factorization gives
\(\operatorname{dist}(z,\Gamma)\ge B(|W(z)|-1)\). Hence
\begin{equation}\label{eq:tube-e}
 0\le e(z)\le\delta/B.
\end{equation}

Let \(q(z)=1-(\Re z-\xi)^2/A^2-(\Im z)^2/B^2\). On \(D\),
\(|\nabla q|\le2/B\). For \(x\in[0,2]\), \(q(x)>1/5\). A point in the inner
\(\delta\)-collar has \(q(z)\le2\delta/B<1/8\). The level set \(q=1/8\) has
distance more than \(1/32\) from \([0,2]\). By domain monotonicity the
interior Green function is at most \(\log160<6\) there. The superharmonic
barrier \(48q\) dominates it on both boundaries of the collar. By the
maximum principle,
\begin{equation}\label{eq:tube-gD}
 g_D(z,x)\le96\delta/B
       \quad(z\in D\cap T_\delta,\ x\in[0,2]).
\end{equation}

Put \(f=A_0-\phi\). Let \(z\) be interior. Sweeping the positive-source
logarithms in \(A_0\) increases their sum by at most \(3\cdot96\delta/B\).
The linear part of \(\phi\) keeps its mean. By \eqref{eq:tube-green}, the
logarithmic-potential part of \(\phi\) has the favorable sign in \(-\phi\).
Thus \(\int f\,d\omega_z-f(z)\le288\delta/B\). Now let \(z=F(w)\) be
exterior. Its mean coordinate is \(\xi+c/\overline w+d/w\), at distance at
most \((A+B)(|w|-1)\) from \(z\). The change in \(A_0\) is exactly
\(-3e(z)\). The remaining terms give the smaller bound
\([1/2+(A+B)\log2]\delta/B<3\delta/B\). Hence, for the swept measure
\(\nu^\#=\int\omega_z\,d\nu(z)\) of mass \(s\),
\[
 \int f\,d\nu^\#\le\int f\,d\nu+288s\delta/B.
\]

Put \(b=\int e\,d\nu\). The changed interaction energies in
\eqref{eq:contour-energy} increase by at most
\(2sb+2(1-s)b+b=3b\le3s\delta/B\). Hence
\[
 \mathcal E(\kappa+\nu^\#,\mu)
  \le\mathcal E(\kappa+\nu,\mu)+(291/B)s\delta
  <\mathcal E(\kappa+\nu,\mu)+300\delta.
\]
This proves \eqref{eq:tube-transfer}, uniformly for every frequency measure
on the full half-line and every allowed pinned/contour mass split.

Let \(z_1,\ldots,z_d\in\Gamma\). Let \(\sigma_{z,\delta}\) be probability
arclength on the circle of radius \(\delta\) about \(z\). Set
\(\nu_\delta=m^{-1}\sum_{a=1}^d\sigma_{z_a,\delta}\). The circle mean
identity \(\int\log|x-z-\delta e^{iu}|\,du/(2\pi)=\log\max(|x-z|,\delta)\)
implies
\begin{equation}\label{eq:tube-circle}
 \frac1{m^2}\sum_{a\ne b}K_0(z_a,z_b)
 \ge I(\nu_\delta)-\frac d{m^2}\log(1/\delta).
\end{equation}
The mixed potentials from \([0,2]\) and \((-\infty,-1/2]\) are harmonic on
these disks. The function \(\phi\) is subharmonic, also at the interior of
its zero segment. So the circular mean of the contour field \(f\) is at
most its central value. Take \(\delta=m^{-2}\) and \(m\ge8\). Then
\eqref{eq:tube-circle} and \eqref{eq:tube-transfer} together cost at most
\(2d\log m+300\) in the unnormalized energy. Coincident particles give a
zero Vandermonde, so no separation hypothesis is needed. A zero of \(F_g\)
at finite \(m\) only decreases the absolute integrand. No two-sided
logarithmic estimate at finite \(m\) is used.

\section{A lower bound for the continuum energy}\label{sec:dual}

For a measure $\sigma$ write
$L_\sigma(z)=\int\log|z-w|\,d\sigma(w)$. Fix an admissible reference
pair $(\eta_0,\mu_0)$ with $\eta_0=\kappa_0+\nu_0$. Its exact cell masses
and contour weights are given in the supplementary certificate files.
Define
\[
\begin{aligned}
Q_{\log}={}&\iint\log|z-w|\,d\eta_0(z)d\eta_0(w)
 +\iint\log|t-s|\,d\mu_0(t)d\mu_0(s)\\
 &+\iint\log|z+t|\,d\eta_0(z)d\mu_0(t),\\
W_\eta={}&A_0-\phi-2L_{\eta_0}-L_{\mu_0^-},\\
W_\mu(t)={}&B_0(t)-\psi(t)-2L_{\mu_0}(t)-L_{\eta_0}(-t).
\end{aligned}
\]
For constants $h_\eta,h_\mu$ satisfying $W_\eta\ge h_\eta$ on $\Gamma$,
put
\[
P=\int_E(h_\eta-W_\eta)_+\,dx,
\qquad T=\int_{1/2}^{\infty}(h_\mu-W_\mu)_+\,dt.
\]

\begin{lemma}[Dual lower bound]\label{lem:dual-lower}
For every admissible pair of measures,
\begin{equation}
\alpha_\Gamma\ge C_0+Q_{\log}+h_\eta+h_\mu-P-T.
\label{eq:duality-bound}
\end{equation}
\end{lemma}

\begin{proof}
Let $\sigma=\eta-\eta_0$ and $\tau=\mu^- -\mu_0^-$; both have total
mass zero. The frequency variation is written on the reflected negative
axis. Reflection preserves its self-energy, so $I(\tau)=I(\mu-\mu_0)$.
The quadratic remainder after linearization is
\[
I(\sigma)+I(\tau)+I(\sigma,\tau)
 =\tfrac12 I(\sigma)+\tfrac12 I(\tau)+\tfrac12 I(\sigma+\tau)\ge0.
\]
A signed measure of mass zero has nonnegative logarithmic energy. For
smooth measures this is the Dirichlet integral identity. Finite-energy
measures follow by approximation. In the frequency first variation, the
change of variables from $x<0$ to $t=-x$ gives $W_\mu(t)$. Expanding the
energy and using the definitions of $W_\eta,W_\mu$ gives
\[
\mathcal E\ge C_0+Q_{\log}+\int W_\eta\,d\eta+\int W_\mu\,d\mu.
\]
The pinned and contour masses sum to one, so the first integral is at
least $h_\eta-P$. By the frequency cap the second integral is at least
$h_\mu-T$. This proves \eqref{eq:duality-bound}.
\end{proof}

The reference has pinned and frequency densities in $[0,1]$ and the
required total masses. On the contour it is a nonnegative symmetric
Fej\'er mixture. Its cell masses and Fej\'er weights are rational. The
checker evaluates the finitely many algebraic Fourier moments by directed
interval enclosures. The supplementary materials give the complete
reference, the potential formulas and closed endpoint evaluations. The
comparison uses no quadrature value.

Take $h_\eta=-669/200$ and $h_\mu=71/25$. The directed-interval
certificate in the supplementary materials checks the bounds on the full
contour, pinned and frequency covers. It also checks the exact reference
integrals:
\[
P\le\frac{6905963574153}{81920000000000},\quad
T_{[1/2,81/10]}\le\frac{38158094390353}{409600000000000},\quad
Q_{\log}>-\frac{438334707}{10^9}.
\]
The rational cell lists and the directed checker are in the supplementary
archive. For the infinite frequency tail we argue by hand. For
$t\ge81/10$,
\[
W_\mu'(t)\ge\log2-\frac1{2(t-1/2)}+\frac3{t+2}
 -\frac2{t-11/2}-\frac1{t-751/5000}.
\]
The right side is increasing and exceeds $1/50$ at $81/10$. The checker
also verifies
\[
W_\mu(81/10)>71/25=h_\mu.
\]
Thus $W_\mu$ is increasing above this endpoint, and
$(h_\mu-W_\mu(t))_+=0$ for every $t\ge81/10$. So the positive-part
integral $T$ has no tail beyond that point. A rational upper bound for
$\log2$ in \eqref{eq:duality-bound} now gives
\begin{equation}
\alpha_\Gamma>\frac{1992909}{10^6}.
\label{eq:alpha-lower}
\end{equation}

\section{A uniform estimate for finite configurations}\label{sec:configuration}

This section passes from the exact integrand of Section~\ref{sec:contour}
to the continuum energy. The estimate holds for every configuration,
including arbitrarily close contour points and arbitrarily large
frequencies. Let \(h=1/m\). For every selected pinned index and frequency
use the cells
\[
 P_i=[i/m,(i+1)/m),\qquad Q_r=[(r-1)/m,r/m),
\]
\begin{equation}\label{eq:config-cells}
 d\kappa=\sum_{i\in S}1_{P_i}dx,\qquad
 d\mu=\sum_{r\in R}1_{Q_r}dt.
\end{equation}
These measures satisfy both density caps exactly. Their masses are \(k/m\)
and one. The first frequency cell starts at \(1/2\), so no boundary cell is
missing. Set \(\nu_\delta=m^{-1}\sum\sigma_{z_a/m,\delta}\) with
\(\delta=m^{-2}\).

The estimates below are elementary and one-sided in the useful direction.
Their constants do not depend on the selected sets or on the largest
frequency.

\subsection{Scalar weights and fixed frequency products}
The entropy bound for a binomial coefficient gives
\[
 \log U_i\le m\phi(i/m),\qquad \log w_r\le m\psi(r/m).
\]
Now \(\psi'(t)=\log(t/(t-1/2))-\log2\) and
\(\int_{1/2}^1(\psi')_+dt=(1/2)\log2\). Summing endpoint-to-cell
differences over any disjoint collection of cells gives
\begin{equation}\label{eq:config-psi}
 m\sum_{r\in R}\psi(r/m)-m^2\int\psi\,d\mu
 \le (m/2)\log2.
\end{equation}
The field \(B_0\) is increasing, which is favorable for its negative sign.
For the fixed lattice denominator, monotone quadrature gives
\begin{equation}\label{eq:config-Bzero}
 0\le mB_0(t)-\sum_{i\in I_m}\log(t+i/m)
                   -\sum_{j\in J_m}\log(t+j/m)
 \le\log(35/3),\qquad t\ge1/2.
\end{equation}
Hence the frequency one-body error is at most
\(m(\log(35/3)+(1/2)\log2)\) on the entire half-line.

\subsection{Pinned denominators}
Take an interval of \(N\) consecutive sites, starting at \(a\) and
containing \(i\). Put \(n_L=i-a\) and \(n_R=a+N-1-i\). The omitted-site
product satisfies
\begin{equation}\label{eq:config-pinned}
 \log(n_L!n_R!)-(N-1)\log m
 -m\int_{a/m}^{(a+N)/m}\log|i/m-y|dy
 \ge\log\frac{m}{n_R+1}.
\end{equation}
For the proof, compare monotone logarithmic sums with the integrals on each
side of \(i\). On the right use \(\log(n_R!)\ge n_R\log n_R-n_R\) and
\(n_R\log(1+1/n_R)\le1\); the case of zero is direct. The \(I_m\)
contribution is nonnegative. The \(J_m\) contribution is at least
\(-\log2\). Also \(\int_E|(\phi-A_0)'|\le\log168\), with contributions at
most \(\log6\), \(2\log2\) and \(\log7\). Hence the pinned one-body cost is
at most \(m\log336\).

\subsection{Real self-cells and the fixed Vandermonde}
Put
\[
 H(x)=\tfrac12x^2\log|x|-\tfrac34x^2\quad(x\ne0),\qquad H(0)=0.
\]
At lattice distance \(d\ge1\), let
\[
 J_d=\int_0^1\int_0^1\log|d+u-v|\,du\,dv,
 \quad e_d=\log d-J_d\ge0.
\]
Since \(H''(x)=\log|x|\), we have \(J_d=H(d+1)+H(d-1)-2H(d)\). Therefore
\(e_1=3/2-2\log2\), and for \(d\ge2\)
\[
 e_d=\sum_{j\ge1}\frac1{j(2j+1)(2j+2)d^{2j}}
 \le\frac1{12d(d-1)}.
\]
By telescoping and Stirling's formula,
\(\sum_{d\ge1}e_d=(1/2)\log(2\pi)-3/4\). Each selected lattice site has at
most two other sites at a given distance. A single self-cell has log
average \(-\log m-3/2\). Hence any \(n\) distinct lattice sites satisfy
\begin{equation}\label{eq:config-self}
 \sum_{a\ne b}\log|x_a-x_b|
 \le m^2\iint\log|x-y|\,d\lambda(x)d\lambda(y)
                        +n(\log m+\log(2\pi)),
\end{equation}
where \(\lambda\) is the union-of-cells measure. This uses no bound on the
range of sites. The pinned and frequency interactions together cost at most
\(2m(\log m+\log(2\pi))\). The fixed \(\Delta(I_m)^2\Delta(J_m)\) costs at
most the same amount, with its leading constant as in
Section~\ref{sec:contour}.

\subsection{Mixed cells}
The exact ellipse clearances are
\[
 d_+=\operatorname{dist}(\Gamma,[0,2])=751/5000,
 \qquad d_-=\operatorname{dist}(\Gamma,(-\infty,-1/2])=1749/5000.
\]
They follow by minimizing the quadratic distance in \(\cos u\). The left
vertex is the closest point to the first interval. The mean cell
displacement is \(h/2\). The pinned--contour interaction costs at most
\(m/d_+\). The contour--frequency interaction costs at most \(m/(2d_-)\),
and the fixed contour denominators at most \(3m/(2d_+)\). The
pinned--frequency cell displacements are centered. By Taylor's inequality,
with \(i/m+r/m\ge3/4\) and \(h\le1/4\), their total loss is at most
\(1/3\). All bounds are uniform in the contour mass. They use no contour
density or separation assumption.

\subsection{Finite contour product}
For every complex \(z\),
\begin{equation}\label{eq:config-Fg}
 \log|F_g(mz)|\le m\phi(z)+b_m,\qquad
 b_m=5(1+2\log2)+4(1+\log(g+1)).
\end{equation}
The proof also covers a zero segment crossing the contour. Write the cell
discrepancy as \(D(w)=\log|w|-\int_0^1\log|w-u|du\). If \(|w|\le2\), then
\(D(w)\le1+2\log2\). If \(|w|>2\), then \(D(w)\le2/|w|\). At most five
indices have \(|\Re(mz)+q|\le2\). On both remaining sides the reciprocal
sums are at most \(4(1+\log(g+1))\). The denominator \(g!\) has the
favorable monotone quadrature sign. This proves \eqref{eq:config-Fg}. A
zero gives \(-\infty\) and is harmless. The circle comparison of
Section~\ref{sec:tube} then costs at most \(2m\log m\) and puts the contour
particles on \(T_{m^{-2}}\).

Insert these estimates into the absolute value of
\eqref{eq:contour-signed}. Discard the favorable factors \(m^{-2m-2d}\) and
\(\ell^k\le1\). This gives
\[
 \log Q_m(S,R,Z)\le
      -m^2\mathcal E(\kappa+\nu_\delta,\mu)+D_{\rm tube}(m),
\]
\[
\begin{split}
 D_{\rm tube}(m)={}&4m(\log m+\log(2\pi))+2m\log m+mb_m+\tfrac13\\
 &+m\left[\log336+\log(35/3)+\tfrac12\log2
                            +\frac5{2d_+}+\frac1{2d_-}\right].
\end{split}
\]
The symmetrization factor is not part of \(Q_m\); it enters at the
integration. Since \(g+1\le m\) for \(m\ge4\), the explicit constants give
\[
 D_{\rm tube}(m)\le10m\log m+c_*m+\tfrac13,
\]
\begin{equation}\label{eq:config-cstar}
 c_*=4\log(2\pi)+9+\tfrac{21}{2}\log2+
         \log336+\log(35/3)+\frac5{2d_+}+\frac1{2d_-}
       <49.979.
\end{equation}
The rational-series checks in the supplementary materials prove this
rounded inequality. By the tube estimate and \eqref{eq:alpha-lower}, for
every \(R\) and every contour configuration,
\begin{equation}\label{eq:config-final}
 \log Q_m\le-\alpha m^2+10m\log m+c_*m+300+\tfrac13,
       \qquad m\ge8,\quad4\mid m.
\end{equation}

\section{The full frequency sum}\label{sec:frequency}

The absolute expansion in Section~\ref{sec:contour} has frequencies
$R\subset\{g+1,g+2,\ldots\}$ with $|R|=m$. As in \eqref{eq:contour-E},
$E_m(H)$ is its absolute envelope restricted by $\max R\le H$. We first
prove
\begin{equation}\label{eq:frequency-cutoff}
 E_m(\infty)\le\left(1+\frac3{160}2^{-m}\right)E_m(20m).
\end{equation}
For fixed pinned and contour variables, the $R$-dependent factor is
\[
 f(R)=
 \frac{\prod_{r\in R}w_r\,\Delta(R)^2
       \prod_{r\in R,z\in Z}|r+z|}
      {\prod_{r\in R,i\in I_m}(r+i)
       \prod_{r\in R,j\in J_m\setminus S}(r+j)}.
\]
Let $r=\max R>20m$. Replace $r$ by the least unused integer
$s\in\{g+1,\ldots,g+m\}$, and call the result $R'$. Then $s\le3m/2$, and for fixed
$R',r$ there are at most $m$ preimages. Since at most two integers have
any given distance from $s$,
\[
 \prod_{u\in R\setminus\{r\}}|s-u|\ge
 d_m:=\left\lfloor\frac{m-1}{2}\right\rfloor!
      \left\lceil\frac{m-1}{2}\right\rceil!
 \ge e^{-1}\left(\frac{m}{2e}\right)^{m-1}.
\]
Using $s+i\le11m/4$, $s+j\le7m/2$, and the ellipse bounds
$\Re z>-m/6$, $|z|<5m/2$ gives
\[
 \frac{f(R)}{f(R')}
 \le\frac{w_r}{w_s}(77/8)^m(189/16)^{m-|S|}
       \frac{m^{2m}}{r^2d_m^2}
 \le12500^m\frac{m^2}{4r^2}w_r.
\]
The last inequality uses $w_s\ge2^{-3m/2}$ and $e<3$. These estimates
include both frequency-dependent minors and all contour--frequency factors.

For a finite cutoff $M\ge20m$, let $Z_M$ be the sum of $f(R)$ over
$R\subset(g,M]$. Put
\[
 b_m=\frac{m}{1600}12500^m\sum_{r>20m}w_r.
\]
The replacement map gives $Z_M-Z_{20m}\le b_mZ_M$. Since
\[
 \sum_{r>g}w_rt^r=\left(\frac{t}{2-t}\right)^{g+1},
\]
evaluation at $t=15/8$ and $200000\cdot8^{20}<15^{20}$ give
\[
\begin{aligned}
 b_m&\le\frac{15m}{1600}
       \left[50000(8/15)^{20}\right]^m\\
    &<\frac3{320}2^{-m}<\frac12.
\end{aligned}
\]
Hence $Z_M\le(1+(3/160)2^{-m})Z_{20m}$. Monotone convergence proves
\eqref{eq:frequency-cutoff}.

There are $2^m$ pinned subsets and at most
$\binom{39m/2}{m}\le(39e/2)^m$ retained frequency subsets. The normalized
contour length is at most $mA$, where $A=12783/10000$, and
$(mA)^d/d!\le e^{mA}$. These counts, \eqref{eq:config-final} and
\eqref{eq:frequency-cutoff} give
\[
\begin{aligned}
 \log|G_m(\beta)|\le{}&
 -\alpha m^2+10m\log m+(c_*+\log39+1+A)m\\
 &+300+\tfrac13+\log\left(1+\frac3{160}2^{-m}\right).
\end{aligned}
\]
Here $c_*<49.979$ and $c_*+\log39+1+A<56$. The constant terms are less
than $301$. This proves \eqref{eq:analytic}.

\section{Clearing denominators}\label{sec:integerization}

This section proves \eqref{eq:integerizer} for all \(m\in4\mathbb N\). For
\(F\in\mathbb Q[T]\), the Gauss valuation \(v_p(F)\) is the minimum
valuation of its coefficients. Set \(v_p(0)=+\infty\).

\subsection{Denominators and a formal identity in the parameter}

Let \(a\) be a variable for now. Put \(c_i=\binom{i+g}{g}\),
\(U_i=c_i a^i\), \(h_i=\sum_{s=1}^g1/(i+s)\), and
\(Q_s(a)=(a/(a-1))^s-1\). In Section~\ref{sec:source} replace \(2^i\) by
\(a^i\) and \(2^s-1\) by \(Q_s(a)\). This gives the rational coefficient
functions of \(D_gM\) in Lemma~\ref{lem:source-bank}. They follow from
Leibniz's rule and

\[
 \begin{aligned}
 L^{(s)}&=(-1)^{s-1}(s-1)!\bigl((a-1)^{-s}-a^{-s}\bigr),\\
 \beta^{(s)}&=(-1)^{s-1}(s-1)!a^{-s}
          \left(L-\sum_{r=1}^{s-1}\frac{Q_r(a)}r\right).
 \end{aligned}
\]

Put \(U=2m-1\) and \(d_U=\operatorname{lcm}(1,\ldots,U)\). The identity
\begin{equation}\label{eq:int-harmonic}
 c_i h_i=\sum_{s=1}^g\frac{(-1)^{s-1}}s\binom{i+g}{g-s}
\end{equation}
is the derivative at zero of the Vandermonde identity
\(\binom{i+g+x}{g}=\sum_{s=0}^g\binom{i+g}{g-s}\binom{x}{s}\). So the
harmonic denominators need only the indices \(s\le g\le U\). The finite sums
have denominators cleared by \(d_U\) in \(C\) and by \(d_U^2\) in \(B\).
The entries of \(B\) have pole order at most \(g\) at \(a=1\). Scaling rows
gives
\begin{equation}\label{eq:int-Pm}
 P_m(a,T):=d_U^{3m}(a-1)^{gm}G_m(a,T)\in\mathbb Z[a,T].
\end{equation}
Here \(G_m(a,T)\) is the determinant with parameter \(a\). The rest of the
paper uses \(G_m(2,T)\). Its finite rational formulas contain no negative
powers of \(a\).

Work formally in \(\mathbb Q[[a]]\), with \(L_1=\sum_{r\ge1}a^r/r\), \(L_2=\sum_{r\ge1}a^r/r^2\),
and

\[
 \gamma_i=\tfrac12\left(h_i^2-\sum_{s=1}^g(i+s)^{-2}\right).
\]

For all \(g,i\ge0\) the rational functions satisfy
\begin{equation}\label{eq:int-t}
 t_i^{(g)}-L_1U_i=\sum_{k\ne i}\frac{c_ka^k}{i-k}-U_i h_i,
\end{equation}
\begin{equation}\label{eq:int-b}
 b_i^{(g)}+L_1U_i h_i+L_2U_i
       =\sum_{k\ne i}\frac{c_ka^k}{(i-k)^2}+\gamma_iU_i.
\end{equation}
Here and below sums over \(k\) start at zero. To prove these uniformly in
\(g\), adjoin to \(\mathbb Q((a))\) a formal symbol \(\lambda\) with
\(\lambda'=1/a\). Differentiating \(D_g a^x=\binom{x+g}{g}a^x\) in the
auxiliary exponent gives

\[
 D_g(a^i\lambda)=U_i(\lambda+h_i),\quad
 D_g(a^i\lambda^2)=U_i(\lambda^2+2h_i\lambda+2\gamma_i).
\]

The substitutes \(\bar L=-L_1-\lambda\) and
\(\bar\beta=-L_2-\lambda^2/2\) obey the differential equations of \(L\) and
\(\beta\): \(\bar L'=1/(a(a-1))\) and \(\bar\beta'=\bar L/a\). The base
formal identities \(t_i^0-L_1a^i=\sum_{k\ne i}a^k/(i-k)\) and
\(b_i^0+L_2a^i=\sum_{k\ne i}a^k/(i-k)^2\) hold coefficientwise. Apply
\(D_g\); it multiplies each \(a^k\) coefficient by \(c_k\). Then compare the
terms free of the independent symbol \(\lambda\). This proves
\eqref{eq:int-t} and \eqref{eq:int-b}, including the coefficient at
\(k=i\). The formula for \(\gamma_i\) also follows from the coefficient of
\(x^2\) in the same Vandermonde identity.

\subsection{Rank-one features and specialization at 2}

For an independent variable \(X\) form

\[
 \mathcal H(a,X)=\begin{bmatrix}C(a)\\
                B(a)-L_1C(a)+L_2A(a)-XA(a)\end{bmatrix}.
\]

Fix a weight \(k\ge0\). On \(I_m\) put \(e_I(i)=[i=k]\) and
\(w_I(i)=v_I(i)^2\), where \(v_I(i)=1/(i-k)\) for \(i\ne k\) and \(v_I\) is
zero otherwise. Define \(e_J,v_J,w_J\) on \(J_m\) in the same way. By
\eqref{eq:int-t} and \eqref{eq:int-b},
\begin{equation}\label{eq:int-rank-one}
 [a^k]\mathcal H=c_k\left(
 \begin{bmatrix}e_I\\v_I-h_ke_I\end{bmatrix}v_J^t
 +\begin{bmatrix}0\\-e_I\end{bmatrix}w_J^t
 +\begin{bmatrix}v_I-h_ke_I\\-w_I-h_kv_I+(\gamma_k-X)e_I\end{bmatrix}e_J^t
                  \right).
\end{equation}
The second rank-one summand vanishes unless \(k\in I_m\), and the third
unless \(k\in J_m\). So weight \(k\) has at most
\(1+[k\in I_m]+[k\in J_m]\) features. A nonzero determinant term must
choose \(2m\) distinct features. For a fixed coefficient of \(a\) only
finitely many features matter. So finite multilinearity proves this
statement formally.

For \(m=4h\), the cheapest \(8h\) features are two copies of
\(0,\ldots,h-1\) and three of \(h,\ldots,3h-1\). Their total weight is
\begin{equation}\label{eq:int-Nm}
 N_m=2\sum_{k=0}^{h-1}k+3\sum_{k=h}^{3h-1}k
                  =13h^2-4h=\frac{13}{16}m^2-m.
\end{equation}
The determinant-preserving row shear gives
\(\det\mathcal H(a,X)=G_m(a,X-L_2)\). Hence the right side is divisible
by \(a^{N_m}\) in \(\mathbb Q[[a]][X]\). The substitution \(X=T+L_2\) is
polynomial in \(X\), so it preserves divisibility. With the integer
polynomial \eqref{eq:int-Pm} this gives
\begin{equation}\label{eq:int-divisible}
 d_U^{3m}(a-1)^{gm}G_m(a,T)\in a^{N_m}\mathbb Z[a,T].
\end{equation}
Only now set \(a=2\). The factor \((a-1)^{gm}\) is one, so
\begin{equation}\label{eq:int-twoadic}
 v_2(G_m)\ge N_m-3m\lfloor\log_2U\rfloor.
\end{equation}
No value of \(L_1(2)\) or \(L_2(2)\) is used. The formal translation is
undone before the integer specialization. The supplementary reports test
finitely many formal coefficients. These tests only corroborate the
argument, which holds for every \(g\).

\subsection{Valuation bounds at odd primes}\label{subsec:odd}

Let \(p\) be an odd prime with \(p^2>U+g\). Put

\[
 r=g\bmod p,\quad \kappa_n=[n\bmod p\ge p-r],\qquad
 \tau_n=(\kappa_n-1)[n\ge p\ \hbox{or}\ p\le g].
\]

When \(r=0\), \(\kappa_n=0\). Since \(n+g<p^2\), Kummer's carry criterion
gives \(v_p\binom{n+g}{g}=\kappa_n\). If \(p\) divides a denominator index
\(s\), then \(v_p\binom{n+g}{g-s}\ge\kappa_n\); if \(s\le n\), then
\(v_p\binom{n-s+g}{g}=\kappa_n\). All nonzero denominator indices and index
differences have valuation at most one. The ultrametric inequality, applied
to the formulas of Section~\ref{sec:source}, gives the following upper
bounds for negative Gauss valuations.

\begin{center}
\begin{tabular}{p{0.36\textwidth}p{0.54\textwidth}}
\hline
Entry & Upper bound \\
\hline
\(C_{ij}\), \(i\ne j\) & \([i\equiv j\pmod p]-\min(\kappa_i,\kappa_j)\) \\
\(C_{ii}\) & \((1-\kappa_i)[p\le g]\) \\
\(B_{ij}\), \(i\ne j\) & \([i\equiv j\pmod p]-\min(\tau_i,\tau_j)\) \\
\(B_{ii}-TU_i\), \(i\ge p\) or \(2p\le g\) & \(2-\kappa_i\) \\
\(B_{ii}-TU_i\), \(i<p\) and \(p\le g<2p\) & \(1\) \\
\(B_{ii}-TU_i\), \(i<p\) and \(p>g\) & \(0\) \\
\hline
\end{tabular}
\end{center}

For the last three lines, the first diagonal sum has a squared \(p\)
denominator only if \(i\ge p\). Both denominators in its double sum are
divisible by \(p\) only if \(2p\le g\), and then the binomial gains
\(\kappa_i\). A single \(p\) denominator costs at most one. The coefficient
\(-TU_i\) has nonnegative valuation and does not worsen the bounds.

Let \(W_m(p)\) be the maximum weight sum over permutations. By Leibniz,
\(-v_p(G_m)\le W_m(p)\). Integer row and column potentials whose sums
dominate every entry weight also give an exponent, namely their total. No
assignment optimality or determinant rank theorem is needed. The coarse
bound is \(W_m(p)\le3m\).

\subsection{Small primes}

Let \(p<g\). After enlarging the special diagonal bounds if needed, take
\(w_C(i,j)=[i\equiv j\pmod p]-\kappa_i\kappa_j\) and \(w_B=1+w_C\). Let
\(c_I(v),c_J(v)\) be the residue copy counts and \(K_I,K_J\) the killed
index counts. Put
\[
 S_U=\sum_{\kappa_v=0}\min(2c_I(v),c_J(v)),\quad
 S_K=\sum_{\kappa_v=1}\min(2c_I(v),c_J(v)).
\]
Then
\begin{equation}\label{eq:int-W}
 W_m(p)\le\min\{m+S_U,\ 3m-2K_I-K_J+S_K\}.
\end{equation}
For the proof, subtract the bottom-row constant \(m\). The score is then
\(x-z\), where \(x\) counts unkilled same-residue matches and \(z\) killed
off-residue matches. Let \(y\) count killed same-residue matches. Then
\(x\le S_U\) and \(y\le S_K\). Counting the killed rows and columns gives
\(z+y\ge2K_I+K_J-2m+x\). Both bounds follow.

For \(\rho=p/m\) and \(t=v/m\in[0,\rho)\), the copy counts are
\[
 c_I(t)=\#\{n:1/4\le t+n\rho<5/4\},\quad
 c_J(t)=\#\{n:0\le t+n\rho<2\}.
\]
They change only at the six endpoints
\begin{equation}\label{eq:int-endpoints}
 0,\rho,\quad a-\lfloor a/\rho\rfloor\rho\ (a=1/4,5/4,2),
 \quad-1/2+(\lfloor1/(2\rho)\rfloor+1)\rho.
\end{equation}
The last is the killed threshold. Every endpoint is \(b+n\rho\) with
integer \(n\) and \(b\in\{0,1/4,5/4,2,-1/2\}\). So all walls lie among the
numbers \(|b-c|/n\), \(n\ge1\). Each fixed positive strip contains finitely
many of them. Within a wall cell the objectives in \eqref{eq:int-W} are
affine. We integrate their lower envelope after also splitting at the
affine crossings. Each endpoint in \eqref{eq:int-endpoints}, multiplied by
\(m\), is an integer. With half-open conventions the residue counts are
therefore exactly \(m\) times the cell widths, with the correct boundary
assignments. No asymptotic equality to the true assignment profile is used.

Let \(\delta=1/128\). The 756 affine pieces in the certificate specify the
exact integral \(L_*\) on \([\delta,1/2]\). Its value is the following
numerator divided by the denominator.
\begin{center}
\begin{minipage}{0.92\linewidth}\footnotesize\ttfamily
\seqsplit{1448276314740194964855260100800462127260298031042796749529725627370010396734895633658285333350450905704443029571002838710237}
\par\medskip /\par\medskip
\seqsplit{1404098881580803157259376918521625735839461675553873436297217652662216162782957043593966060046854826150978176673015236224000}
\end{minipage}
\end{center}

Near zero the row potentials \(1-\kappa_i\) and \(2-\kappa_i\), with zero
column potentials, are feasible. They give \(W_m(p)\le3m-2K_I\). Any
interval of \(m\) integers contains at least
\(\lfloor m/p\rfloor(g\bmod p)\) killed indices. Hence
\begin{equation}\label{eq:int-near-zero}
 W_m(p)/m\le3-2\{1/(2\rho)\}+2\rho.
\end{equation}
Write \(\delta=1/(2Q)\), \(Q=64\). The substitution \(x=1/(2\rho)\) gives
\[
 \int_0^\delta(3-2\{1/(2\rho)\})d\rho
 =\frac1{2Q}+\sum_{q\ge Q}\left(\frac1q-\log(1+1/q)\right).
\]
The summand exceeds \(1/(2q(q+1))\) by exactly
\(\int_0^1s(1-s)/(q(q+s)(q+1))\,ds\le1/(6q^2(q+1))\). Summing, and adding
the integral of \(2\rho\), gives the rigorous upper bound
\begin{equation}\label{eq:int-small}
 I_{\rm small}\le2\delta+2\delta^2=129/8192.
\end{equation}

\subsection{Large primes and the integerizer}

For \(\rho>1/2\), keep each row copy in both blocks and every column copy.
In a residue cell with flag \(\kappa\), copy digit \(t\) has
\(\tau=(\kappa-1)[t\ge1]\). The exact data give an integer potential to
each row and column type. Every pair of types must satisfy the off-residue
bound of Section~\ref{subsec:odd}. A pair in the same residue cell must also
satisfy the same-residue bound, including the diagonal case. Checking both
cases also covers two distinct residues in the same cell. The objective is
the sum of potentials times the exact cell counts.

The exact rational verifier in the supplementary materials reconstructs
the complete wall geometry. It checks all 756 low pieces and all 1,942 high
integer inequalities. It gives the following high upper profile.

\begin{center}
\begin{tabular}{p{0.36\textwidth}p{0.54\textwidth}}
\hline
Interval in \(\rho\) & Upper profile \\
\hline
\((1/2,7/12)\) & \(1\) \\
\((7/12,5/8)\) & \(3\rho -3/4\) \\
\((5/8,2/3)\) & \(\rho +1/2\) \\
\((2/3,3/4)\) & \(4\rho -3/2\) \\
\((3/4,5/6)\) & \(3\rho -3/4\) \\
\((5/6,1)\) & \(7/4\) \\
\((1,7/4)\) & \(15/4-2\rho\) \\
\((7/4,2)\) & \(2-\rho\) \\
\((2,5/2)\) & \(0\) \\
\hline
\end{tabular}
\end{center}

Its integral is \(287/192\). The copy geometry splits this table into 11
cells. The data list the individual types and integer potentials. No
optimizer or numerical rounding is used. On each wall in a positive strip
use \(3m\) instead. This costs \(O(m\log m)\) in total and covers
\(\rho=1/2\).

Choose the exponents \(E_p\) in the following order.
\begin{enumerate}
\item At 2 take \(E_2=3m\lfloor\log_2U\rfloor-N_m\).
\item At odd \(p^2\le U+g\) use \(3m\lfloor\log_p U\rfloor\).
\item At \(p>U\) use zero.
\item At the remaining \(p/m<\delta\) use \(3m-2K_I\).
\item At a wall in a positive strip use \(3m\).
\item Elsewhere below \(m/2\) use \eqref{eq:int-W}.
\item Elsewhere above \(m/2\) use the least feasible high dual objective.
\end{enumerate}
The positive rational product
\begin{equation}\label{eq:int-Dm}
 D_m=\prod_p p^{E_p}
\end{equation}
is finite, and every exponent bounds \(-v_p(G_m)\). So \(D_m\) integerizes
every coefficient. A negative \(E_2\) only divides out a proved common power
of two.

The positive 2-adic overhead is \(O(m\log m)\). Small odd primes cost at
most \(3m\,\pi(\sqrt{U+g})\log U=O(m^{3/2}\log m)=o(m^2)\). On each fixed
positive strip we use the classical weighted prime number theorem
\[
 \frac1m\sum_p f(p/m)\log p\longrightarrow\int f(\rho)d\rho
\]
for a bounded finitely piecewise smooth function \(f\). For
\eqref{eq:int-near-zero}, first omit \(p<\varepsilon m\). Then apply the
theorem on \([\varepsilon,\delta]\), and let \(\varepsilon\downarrow0\). By
\(3m\theta(\varepsilon m)\), the omitted part has normalized limsup at most
\(3\varepsilon\). The exact counts on positive strips and the finite wall
cost now prove

\[
 \limsup\frac{\log D_m}{m^2}
 \le L_*+\frac{129}{8192}+\frac{287}{192}-\frac{13}{16}\log2.
\]

Finally let \(\ell_-=2\sum_{k=0}^{29}(1/3)^{2k+1}/(2k+1)<\log2\). The same
verifier checks the exact rational comparison

\[
 \begin{aligned}
 \tau_{\rm basic}:={}&L_*+\frac{129}{8192}+\frac{287}{192}
                            -\frac{13}{16}\ell_-\\
 ={}&1.97881984510269917464907037626379\ldots<1.979.
 \end{aligned}
\]

The decimal is the exact rational of the receipt, rounded. This proves
\eqref{eq:integerizer}, and in particular the weaker rate \(\tau\le1.99\).

\section{Nonvanishing on a prime-power sequence}\label{sec:nonvanishing}

Fix \(M\in4\mathbb N\), set \(h=M/2\), and choose an odd prime \(p>2M+h\).
For every \(k\ge0\) put \(q=p^k\), \(m=Mq\), \(g=hq\). The scaled pencil
satisfies
\begin{equation}\label{eq:nv-pencil}
\begin{aligned}
 \mathcal H_k(X)&:=q^{3m}G_m(X/q^2)
 =\det\begin{bmatrix}qC^{(g)}\\q^2B^{(g)}-XA^{(g)}\end{bmatrix}
 \in\mathbb Z_{(p)}[X],\\
 \overline{\mathcal H_k}(X)&=\overline{G_M}(X)^q.
\end{aligned}
\end{equation}
Every nonzero denominator in the formulas of Section~\ref{sec:source} has
absolute value less than \((2M+h)q<pq\). So \(q/d\) is \(p\)-integral. It
reduces to zero unless \(q\mid d\), and to \(e^{-1}\) at \(d=eq\),
\(0<|e|<p\). Likewise \(q^2/(d_1d_2)\) vanishes on reduction unless both
denominators are divisible by \(q\).

Write \(i=aq+v\), \(j=bq+w\), \(0\le v,w<q\). Lucas and Frobenius give
\begin{equation}\label{eq:nv-lucas}
 \binom{(a+h)q+v}{hq}\equiv\binom{a+h}{h},\qquad
 U_{aq+v}^{(hq)}\equiv2^v U_a^{(h)}\pmod p.
\end{equation}
Multiply by \(q\). In the first sum in \(t_i\), only the denominator
indices \(dq\) with \(1\le d\le a\) survive. In the second, only \(dq\) with
\(1\le d\le h\) survive. Since \(q\) is odd, parity is preserved. Also
\(2^{dq}-1\equiv2^d-1\pmod p\), and Lucas reduces the binomial to the seed
one. Therefore \(qt_{aq+v}^{(hq)}\equiv2^vt_a^{(h)}\pmod p\). If \(v\ne
w\), then \(q\nmid j-i\), so both \(qC_{ij}\) and \(q^2B_{ij}\) reduce to
zero. If \(v=w\) and \(a\ne b\), the divided differences reduce to \(2^v\)
times their seed entries.

On the diagonal, the multiples of \(q\) in \([i+1,i+hq]\) are exactly
\((a+1)q,\ldots,(a+h)q\). This holds even when \(v=0\) or \(q-1\). Thus
\(qC_{ii}\equiv2^vC_{aa}^{(h)}\). In the first diagonal \(B\) sum only the
squared denominators \((dq)^2\) survive. In the double sum both indices must
be multiples of \(q\), with \(2\le d\le h\), \(1\le e<d\). Lucas, parity and
Frobenius reduce these terms to the seed double sum. Hence
\(q^2B_{ii}\equiv2^vB_{aa}^{(h)}\) and \(A_{ii}\equiv2^vA_{aa}^{(h)}\). All
scaled entries are integral before reduction. The argument holds for each
\(k\) directly, without a recurrence in \(k\).

The window endpoints are multiples of \(q\). Grouping rows and columns by
residue gives, in each residue block, the seed pencil times \(2^v\). The
row and column permutations are the same transpose permutation of a
\(2M\times q\) array, so their signs cancel. The product of the scalar
block determinants is \(2^{Mq(q-1)}\equiv1\pmod p\), because
\(p-1\mid q-1\). This proves \eqref{eq:nv-pencil}, including the exact
scalar and sign.

The modular data in the supplementary materials show that for \(M=32\),
\(p=83\) the matrix of Section~\ref{sec:source} satisfies
\begin{equation}\label{eq:nv-seed}
 \det\begin{bmatrix}C^{(16)}\\B^{(16)}\end{bmatrix}
                           \equiv30\pmod{83}.
\end{equation}
This is a finite \(64\times64\) matrix check. The saved entries, row swaps
and unnormalized Gaussian pivots determine the determinant by integer
modular arithmetic. The code also rebuilds the matrix with exact rationals,
and independently by a modular Pascal recurrence. No larger determinant is
computed.

Let \(M=32,p=83\), and fix \(t\in\mathbb Q_{83}\) with
\(2k+v_{83}(t)>0\). Then \(X=q^2t\) lies in \(83\mathbb Z_{83}\). By
\eqref{eq:nv-pencil} and \eqref{eq:nv-seed},
\begin{equation}\label{eq:nv-final}
 q^{3m}G_m(t)=\mathcal H_k(q^2t)\equiv30^q=30\pmod{83},
 \qquad v_{83}(G_{32\cdot83^k}(t))=-3k(32\cdot83^k).
\end{equation}
This proves \eqref{eq:nonvanishing}. For \(t=a/b\) in lowest terms,
\(k\ge\lfloor v_{83}(b)/2\rfloor+1\) suffices; \(v_p(0)=+\infty\) covers
zero. The sequence does not depend on \(t\), and every index lies in
\(32\mathbb N\).


\section{Proof of Theorem~\ref{thm:main}}\label{sec:proof}

Suppose $\beta=a/b$ with coprime integers $a,b$ and $b>0$. Along
$m=32\cdot83^k$, the condition in \eqref{eq:nonvanishing} holds for all large
$k$. Since $\deg G_m\le m$ and $D_mG_m\in\mathbb Z[T]$,
\[
 Z_m=b^mD_mG_m(a/b)\in\mathbb Z\setminus\{0\}.
\]
For large $m$, \eqref{eq:integerizer} gives $\log D_m\le1.979m^2$. By
\eqref{eq:analytic},
\[
 \log|Z_m|\le -0.013909m^2+10m\log m+(56+\log b)m+301
 \longrightarrow -\infty.
\]
This contradicts $|Z_m|\ge1$.
\qed

\appendix

\section{Computer checks}\label{app:ver}

The proof uses three finite computations.
\begin{enumerate}
\item Directed rational-interval bounds for the dual certificate of
  Section~\ref{sec:dual}.
\item Exact rational reconstruction of the residue cells and integer
  potentials of Section~\ref{sec:integerization}.
\item The determinant of the $64\times64$ seed matrix modulo $83$ in
  Section~\ref{sec:nonvanishing}.
\end{enumerate}
The supplementary archive contains the finite data, the checking programs,
the environment details and the replay outputs. Its SHA-256 and the replay
commands are listed in \texttt{SUPPLEMENTARY.md}. The full replay takes about
one minute on one core of an aarch64 Linux server.

\section*{Use of AI tools}

The author devised the method of proof. Codex (OpenAI, GPT-6 Astra and GPT-6 Luna) searched for the parameters in \eqref{eq:main-polynomial}
and wrote the Python code and the certificates. Claude Code (Anthropic,
Claude Opus 5.5) checked the code, the parameters and the certificates. The author
checked every step of the proof, every table and every reference, and takes
full responsibility for the content of the paper.

\section*{Statements and declarations}
\textbf{Competing interests.} The author has no competing interests.

\textbf{Funding.} No funding was received for this work.

\textbf{Data availability.} The exact data and the verification scripts are
in the supplementary material of this paper.

\begin{thebibliography}{10}

\bibitem{CDT}
F.~Calegari, V.~Dimitrov, and Y.~Tang, \emph{The linear independence of $1$,
$\zeta(2)$, and $L(2,\chi_{-3})$}, preprint (2024), arXiv:2408.15403v2.

\bibitem{DLMF}
\emph{NIST Digital Library of Mathematical Functions},
\url{https://dlmf.nist.gov/}, Release 1.2.8 of 2026-09-15, F.~W.~J. Olver
et al., eds.

\bibitem{Ha}
M.~Hata, \emph{Rational approximations to the dilogarithm}, Trans. Amer. Math.
Soc. \textbf{336} (1993), no.~1, 363--387,
\url{https://doi.org/10.1090/S0002-9947-1993-1147401-5}. MR1147401

\bibitem{Pr}
T.~Prellberg, \emph{Three irrationality results for the dilogarithm:
Rhin--Viola and Viola--Zudilin constructions at $-1/4$, $1/5$, and $-1/3$},
preprint (2026), arXiv:2609.06895v1.

\bibitem{RV96}
G.~Rhin and C.~Viola, \emph{On a permutation group related to $\zeta(2)$},
Acta Arith. \textbf{77} (1996), no.~1, 23--56,
\url{https://doi.org/10.4064/aa-77-1-23-56}. MR1404975

\bibitem{RV05}
G.~Rhin and C.~Viola, \emph{The permutation group method for the dilogarithm},
Ann. Sc. Norm. Super. Pisa Cl. Sci. (5) \textbf{4} (2005), no.~3, 389--437,
\url{https://doi.org/10.2422/2036-2145.2005.3.02}. MR2185863

\bibitem{RV19}
G.~Rhin and C.~Viola, \emph{Linear independence of $1$, $\mathrm{Li}_1$ and
$\mathrm{Li}_2$}, Mosc. J. Comb. Number Theory \textbf{8} (2019), no.~1,
81--96,
\url{https://doi.org/10.2140/moscow.2019.8.81}. MR3864310

\bibitem{VZ}
C.~Viola and W.~Zudilin, \emph{Linear independence of dilogarithmic values},
J. Reine Angew. Math. \textbf{736} (2018), 193--223,
\url{https://doi.org/10.1515/crelle-2015-0030}; preprint version:
Max-Planck-Institut f\"ur Mathematik Preprint Series 2014 (26),
\url{https://archive.mpim-bonn.mpg.de/id/eprint/1914/}. MR3769989

\end{thebibliography}
\end{document}
